% Realiza gradient descent en una funcion modelo con contornos elipticos
% Necesita que este definido:
%  - path(c)   : el mínimo
%  - path(x0)  : el initial guess
%  - \b        : aspect ratio de la elipse
%  - \learning : factor de aprendizaje constante

% Obtengo el centro
\path(c); \getlastxy
\pgfmathsetmacro{\cx}{\x}
\pgfmathsetmacro{\cy}{\y}

\fill[celeste](x0) circle(3pt) node[above right] {$x_0$};

\foreach \m in {1,2,...,8} {

  % Posicion
  \path(x0); \getlastxy

  % Gradiente
  \pgfmathsetmacro{\dx}{2.*(\x-\cx)} % 2(x-x0)/a^2
  \pgfmathsetmacro{\dy}{2.*(\y-\cy)/(\b*\b)} % 2(y-y0)/b^2
  \pgfmathsetmacro{\norm}{sqrt(\dx*\dx+\dy*\dy)}
  \pgfmathsetmacro{\vx}{-\learning*\dx/\norm}
  \pgfmathsetmacro{\vy}{-\learning*\dy/\norm}

  % Draw gradiente
  \draw[->,shorten >=5, thick, amarillo] (x0) --  ++(\vx,\vy);

  \path(x0) ++(\vx,\vy) coordinate(x\m);
  \fill[celeste](x\m) circle(3pt) node[above right] {$x_\m$};

  % Save last position as the new x0
  \path(x\m) coordinate(x0);
}

 
% Contornos de la funcion  
\foreach \i in {1,2,...,9} {
\pgfmathsetmacro\lab{\i+0.1}
\scriptsize\draw[rosa,thick] (c) ellipse  [x radius=\i*1cm, y radius=\i*\b*1cm] node[xshift=\lab cm]{\i};
}
     
